% Authoritative LaTeX manuscript; Markdown is a generated compatibility projection.

\section{Problem}

A document-processing service can exhaust worker memory even when its individual transformations are correct. Dividing work by page may reduce intermediate allocations while leaving large media payloads resident. Increasing concurrency can then improve latency on small inputs and worsen failure on large ones. The engineering problem is to choose a representation and execution policy that preserve resource bounds while allowing useful parallel work.

\section{Conceptual Design}

% Evidence: E26 | source: book/data/local-evidence/evidentiary-basis-11202026/integration/records/E26.md | kind: dated-case-evidence
DocAble's media-separation design distinguishes document structure from aggregate media residency. A streaming transformation extracts media into external storage and retains references in a smaller structural skeleton. Operations that need only the structure can use the skeleton; media consumers resolve the references or require a representation with the media restored. Recombination must also stream: buffering the complete payload on output recreates the memory problem. \autoref{interlude-execution-unit} contrasts this arrangement with whole-document and page-partitioned processing.

Including baseline resident memory, the whole-worker residency requirement uses the same bound as \autoref{4.2-one-problem-many-models:heading:remove-media-from-the-working-set}:
\[
\begin{aligned}
\operatorname{RSS}_{\mathrm{peak}} &\le C,\\
C &= M_{\mathrm{base}} + M_{\mathrm{stream}} + M_{\mathrm{skeleton}}
   + M_{\mathrm{handles}} + M_{\mathrm{work}}.
\end{aligned}
\]
The baseline term \(M_{\mathrm{base}}\) bounds the runtime and other baseline resident memory, including memory-management headroom. The remaining terms bound the rolling stream buffers, resident structural skeleton, handles, and active remediation working set. Within the declared operating envelope, \(C\) is independent of aggregate media payload size. Architecture bounds stream buffers; workload limits bound structural size and reference count; configuration limits active concurrency. Measurements must establish the working set of each media consumer, including decoded images, rendering, and temporary allocations. Externalizing bytes alone does not establish the bound.

Dependency declarations then identify ready computations, while mutation restrictions and resource limits constrain which may execute together. \autoref{chapter:appendix-a-remediation-computation-graph} describes those declarations and their consumers. A performance model evaluates candidate policies using operation costs, dependency paths, concurrency limits, and deployment overhead.

The representations and execution decisions fit together as shown in \autoref{handbook-resource-architecture}.

\begin{figure}[htbp]
\centering
\includegraphics[width=\linewidth]{book/latex/assets/handbook-resource-architecture-062243c02e.pdf}
% Evidence: E26 | source: book/data/local-evidence/evidentiary-basis-11202026/integration/records/E26.md | kind: implementation
\caption[Resource-aware remediation execution]{\emph{Resource-aware remediation execution.} Streaming extraction separates document structure from media payloads. The computation graph determines which work is ready; configured concurrency limits, calibrated against the memory envelope, restrict active work. Media is retrieved on demand and restored through streamed recombination. Residency and performance models guide these decisions, while measurements test their assumptions.}
\label{handbook-resource-architecture}
\end{figure}

\section{Engineering Models}

\autoref{handbook-resource-models} identifies the representations needed to connect resource feasibility to execution design.

% Evidence: E26 | source: book/data/local-evidence/evidentiary-basis-11202026/integration/records/E26.md | kind: dated-case-evidence
% Evidence: E27 | source: book/data/local-evidence/evidentiary-basis-11202026/integration/records/E27.md | kind: measurement-and-analysis
% Evidence: E28 | source: book/data/local-evidence/evidentiary-basis-11202026/integration/records/E28.md | kind: derived-analysis
{\footnotesize
\begin{longtable}{@{}p{.20\linewidth}p{.36\linewidth}p{.36\linewidth}@{}}
\caption[Engineering models for resource-aware execution]{\emph{Engineering models for resource-aware execution.} Representation, scheduling, and measurement must agree on the work and the resource boundary being evaluated.}\label{handbook-resource-models}\\
\toprule
\textbf{Model} & \textbf{Meaning and authority} & \textbf{Consumers} \\
\midrule\endhead
Residency & Retained objects, allocation sizes, release conditions, and an admissible memory envelope. & Capacity decisions, concurrency feasibility, and resource guards. \\
Document structure & Skeleton entities and references to media that need not remain resident. & Remediation passes, media access, and streamed recombination. \\
Execution & Computation prerequisites, readiness, mutation restrictions, and allowed concurrency. & Scheduling policies and execution-order tests. \\
Performance & Operation costs, critical paths, fan-out costs, and deployment assumptions. & Predictions and comparisons of candidate designs. \\
Measurement & Workload identities, measured quantities, sampling boundary, and configuration. & Comparisons between predicted and observed memory, latency, and cost. \\
\bottomrule
\end{longtable}
}

The document format determines the identities in the structural model. The implementation supplies the actual consumers and allocation behavior; the engineer supplies the admissible operating envelope. Measurement records must distinguish managed allocation, process RSS, and container memory. A bound on one does not automatically bound the others, particularly when renderers run in child processes or temporary files occupy memory-backed storage.

\section{Realization}

\subsection*{Preserve Structure and Resolve Media}

% Evidence: E26 | source: book/data/local-evidence/evidentiary-basis-11202026/integration/blueprint-A9-archaeology.md | kind: implementation
For OOXML, the skeleton retains non-media package entries, including document XML and relationships. A manifest records each reference site separately from each distinct payload. Its essential declaration is:

\begin{verbatim}
MediaManifestEntry(
    Anchor, Key, OriginalPartName, RestoreState?)
StoredMediaBlob(ContentHash, FileName, ByteSize)
\end{verbatim}

An OOXML anchor identifies a package path, such as \texttt{ppt/media/image1.png}, rather than a slide ordinal. The key is the SHA-256 hash of the payload. Several reference sites can therefore resolve to one stored blob while retaining their own structural identities. Reordering slides must not change which media each site references.

Extraction copies each payload through a bounded buffer while computing its hash, then stores the file under that hash. The pipeline uploads each distinct file and retains a hash-to-handle map. Recombination copies the revised structural package and restores each referenced media payload to its original package location through bounded streaming. Missing references are errors: an apparently successful output must not silently omit media.

PDF requires a different representation. A sentinel inside an image XObject carries the content-hash identity through intervening document rewrites. The manifest also preserves the stream's filter and decoding parameters. Recombination restores those dictionary entries with the original bytes and flushes each object. A package path cannot supply this identity, and a PDF object number is not stable enough across rewrites.

% Evidence: E27 | source: book/data/local-evidence/evidentiary-basis-11202026/integration/records/E27.md | kind: measurement-and-analysis
Check the consumer before placing extraction in the pipeline. DocAble's PDF alt-text and presentation-composition operations consume whole-page or whole-slide renders. Removing embedded images before rendering leaves blank regions and destroys evidence the reasoner needs. Media separation is consequently useful only where the required visual content remains available; it does not justify sending every consumer a stripped skeleton. Rendering also introduces a decoded working set that the encoded payload size does not determine.

\subsection*{Test the Bound and Its Explanation}

% Evidence: E26 | source: book/data/local-evidence/evidentiary-basis-11202026/integration/blueprint-A9-archaeology.md | kind: test-definition
The streaming arrangement has two independent obligations: preserve the document's content and keep memory consumption within the intended envelope. Restoration tests examine byte identity, structural edits, deduplication, stable references after reordering, and missing-media failures. Resource tests vary payload size and compare the streaming operation with a deliberately buffered alternative. The latter supplies a negative control: if the measurement cannot distinguish the buffered implementation's larger allocations, it is not adequate evidence for the claimed resource property.

These tests establish properties of the selected operation, not the entire process. DocAble's Python streaming test measures allocated memory with \texttt{tracemalloc}, not process RSS. A production resource envelope additionally requires measurements across structural size, decoded image area, concurrent operations, and any child processes involved. In a PDF recombination experiment, per-object flushing keeps peak RSS approximately flat as image count increases; without flushing, memory grows substantially with each additional image. This supports the tested recombination path, not a bound for the whole remediation pipeline.

% Evidence: E27 | source: book/data/local-evidence/evidentiary-basis-11202026/integration/records/E27.md | kind: measurement-and-analysis
A PDF investigation illustrates why the explanation itself needs testing. A working hypothesis attributes gigabytes of memory to a resident document graph in the canonical PDF library. A probe walks an image-heavy PDF and forces collection, revealing a small live graph; releasing every stream reclaims nothing. Memory mapping, capture buffers, serialization, and transient allocations explain the discrepancy. The engineering response changes which holders the model represents and which operations deserve optimization. Increasing the supposed document-graph allowance would preserve the wrong explanation.

\subsection*{Separate Readiness from Resource Admission}

% Evidence: E26 | source: book/data/local-evidence/evidentiary-basis-11202026/integration/blueprint-A9-archaeology.md | kind: implementation-and-model
DocAble's chunkless dependency scheduler waits for each predecessor's result to be applied before treating the consumer as ready. It then acquires a concurrency permit before executing the ready computation. The permit is released when execution finishes, including after failure. Dependency readiness and resource admission are therefore separate decisions.

The configured concurrency limit must fit the resource envelope. The distribution model expresses the active footprint as concurrency limit times per-call working set, independently of the number of waiting units. The semaphore enforces a count, not a byte estimate. The engineer must calibrate the limit for the workload and deployment, or add resource-aware admission where per-call demands vary too much for a count limit.

\subsection*{Predict Critical Paths and Repeated Work}

% Evidence: E28 | source: book/data/local-evidence/evidentiary-basis-11202026/integration/records/E28.md | kind: derived-analysis
Weight each dependency node with an operation duration and compute its earliest finish from the largest predecessor finish. The maximum finish supplies a latency lower bound for the modeled dependencies and durations under sufficient parallel capacity. DocAble's analysis combines measured and estimated durations; resource limits, coordination, and inaccurate estimates can increase observed execution time. The calculation supports comparisons among designs, not a guaranteed end-to-end latency.

% Evidence: E29 | source: book/data/local-evidence/evidentiary-basis-11202026/integration/blueprint-A9-archaeology.md | kind: model
Repeated production adds a different cost. DocAble's time-complexity model uses
\[
T(d) \approx (R+1)F(d)+G(d),
\]
where \(R\) counts production rounds, \(F(d)\) is elapsed time for parsing, traversal, document passes, rendering, and fixed overhead per invocation, and \(G(d)\) is the contribution of reasoning operations and their distribution to elapsed time. The extra invocation applies the results. This approximation assumes that invocation costs are sufficiently similar, that the two contributions do not overlap, and that resolved reasoning units are reused rather than repeated. It explains why shortening a reasoning path can leave repeated capture or rendering dominant. Measure rounds and operation counts as well as elapsed time, and revise the decomposition when the observations disagree.

\subsection*{Retain Structural Performance Checks}

% Evidence: E29 | source: book/data/local-evidence/evidentiary-basis-11202026/integration/records/E29.md | kind: implementation-and-measurement
DocAble's performance ratchet counts production rounds, logical GenAI calls, capture operations, and renders. Small deterministic fixtures establish baselines for these quantities. The comparison rejects execution-mode mismatches; a change that increases a protected count fails the corresponding check unless the baseline is deliberately revised. Repeated runs also test whether the counts are stable.

% Evidence: E29 | source: book/data/local-evidence/evidentiary-basis-11202026/integration/blueprint-A9-archaeology.md | kind: implementation
A fifth quantity, the number of GenAI hops, carries an audit-only label and is excluded from the four upper-bound comparisons. A separate test nevertheless compares it with the seeded value and can fail. The label alone therefore does not establish whether a measurement is advisory; the consumer and its acceptance rule determine the consequence.

Structural counts expose additional work without depending on machine-specific wall-clock measurements. They do not establish production latency or memory consumption, and their coverage is limited to the exercised workflows. Improvements should lower the relevant baseline; increases require an engineering decision about whether the additional work is justified.

\section{Adaptation and Limits}

This arrangement applies when large payloads can be separated from the structure needed for transformation. Its effectiveness depends on preserving reference identity, providing media to consumers that require it, and keeping both extraction and recombination within the intended memory envelope. Consumers that require complete objects, unbounded structural growth, or large decoded working sets may need a different representation or stricter concurrency limits.

The resource bound is conditional on the workload and measurement boundary. Streaming alone does not establish bounded process memory, and a ready computation is not necessarily admissible under the available resource budget. Changes to document formats, rendering libraries, deployment tiers, or concurrency policies require renewed validation of these assumptions.
